\documentclass[12pt]{article}

\usepackage[utf8]{inputenc}
\usepackage[russian]{babel}
\usepackage{amsmath}
\usepackage{amsfonts}
\usepackage{amssymb}

\title{Пример документа LaTeX}
\author{Автор}
\date{\today}

\begin{document}

\maketitle

\section{Введение}
Данная статья представляет собой пример документа, созданного с использованием системы подготовки научных публикаций LaTeX. В этом документе демонстрируется структура стандартной статьи, включающей введение, основную часть и заключение. LaTeX предоставляет мощные средства для форматирования научных и технических документов, обеспечивая высокое качество типографики.

\section{Основная часть}
LaTeX особенно популярен в академической среде благодаря своей способности корректно обрабатывать математические формулы, библиографические ссылки и сложную структуру документов. Ниже приведён пример математической формулы:

\begin{equation}
    E = mc^2
\end{equation}

\section{Заключение}
В заключение можно отметить, что LaTeX остаётся одним из наиболее эффективных инструментов для подготовки научных публикаций. Его использование позволяет сосредоточиться на содержании документа, не отвлекаясь на форматирование, поскольку большинство решений по оформлению определяются заранее установленными стилями. Этот документ успешно демонстрирует базовую структуру статьи с введением и заключением.

\end{document}